\documentclass[12pt]{article}
\usepackage[utf8]{inputenc}
\usepackage[T1]{fontenc}
\usepackage{hyperref}
\title{Suba-Simbiti Language Resources for K-12 Curriculum Development}
\author{Compiled by Claude Sonnet 5.5 (claude-sonnet-5-5) for Andr\'as Kornai}
\date{}
\begin{document}
\maketitle

\section{Introduction}
Suba-Simbiti (ISO 639-3 ssc; Kisuba, Kisimbiti) is a Bantu language spoken in the Mara Region of Tanzania, mainly around Tarime and Musoma. The groups that use it include the Hacha, Kine, Sweta, Simbiti and Kiroba. About 110,000 speakers were reported in 2011. Simbiti is the most widely understood dialect and was chosen for Bible translation. Other listed varieties include Iryege, Kironi, Kiseru, Kisingiri and Rieri. It is closely related to, but distinct from, the Suba language of Kenya (sxb). Do not confuse the two.

\section{Primary Educational Resources}
No primers, textbooks, or school materials were found in search. No evidence of use in schools was found.

\section{Linguistic and Reference Materials}
\begin{itemize}
\item Wikipedia article on Suba-Simbiti, with a summary of its phonology and noun classes: \url{https://en.wikipedia.org/wiki/Suba-Simbiti_language}
\item Wikipedia article on the Suba people of Tanzania: \url{https://en.wikipedia.org/wiki/Suba_people_(Tanzania)}
\item Glottolog note on the split of the older Suba into Kenyan Suba and Tanzanian Suba-Simbiti: \url{https://glottolog.org/resource/languoid/id/suba1254}
\end{itemize}

\section{Religious and Community Resources}
Joshua Project lists the people groups and their districts: \url{https://joshuaproject.net/people_groups/18960}. The Jesus Film has been produced in the language, according to the search summary.

\section{Digital and Audio Resources}
Global Recordings Network has audio material in Suba-Simbiti: \url{https://www.globalrecordings.net/en/language/ssc}. It also has a recording in the Iryege variety: \url{https://www.globalrecordings.net/en/language/25474}.

\section{Curriculum Development Context}
Documentation is thin. Community-led orthography and primer development would have to come first. Tanzanian schooling uses Swahili and English, so a bilingual transition model would be needed.

\end{document}